\documentclass[conference]{IEEEtran}
\IEEEoverridecommandlockouts
% The preceding line is only needed to identify funding in the first footnote. If that is unneeded, please comment it out.
\usepackage{cite}
\usepackage{amsmath,amssymb,amsfonts}
% \usepackage{algorithmic}
\usepackage{algorithm}
\usepackage{algpseudocode}
\usepackage{graphicx}
\usepackage{textcomp}
\usepackage{xcolor}
\def\BibTeX{{\rm B\kern-.05em{\sc i\kern-.025em b}\kern-.08em
    T\kern-.1667em\lower.7ex\hbox{E}\kern-.125emX}}
\begin{document}

\title{Conference Paper Title*\\
{\footnotesize \textsuperscript{*}Note: Sub-titles are not captured in Xplore and
should not be used}
\thanks{Identify applicable funding agency here. If none, delete this.}
}

\author{\IEEEauthorblockN{1\textsuperscript{st} Given Name Surname}
\IEEEauthorblockA{\textit{dept. name of organization (of Aff.)} \\
\textit{name of organization (of Aff.)}\\
City, Country \\
email address or ORCID}
\and
\IEEEauthorblockN{2\textsuperscript{nd} Given Name Surname}
\IEEEauthorblockA{\textit{dept. name of organization (of Aff.)} \\
\textit{name of organization (of Aff.)}\\
City, Country \\
email address or ORCID}
\and
\IEEEauthorblockN{3\textsuperscript{rd} Given Name Surname}
\IEEEauthorblockA{\textit{dept. name of organization (of Aff.)} \\
\textit{name of organization (of Aff.)}\\
City, Country \\
email address or ORCID}
\and
\IEEEauthorblockN{4\textsuperscript{th} Given Name Surname}
\IEEEauthorblockA{\textit{dept. name of organization (of Aff.)} \\
\textit{name of organization (of Aff.)}\\
City, Country \\
email address or ORCID}
\and
\IEEEauthorblockN{5\textsuperscript{th} Given Name Surname}
\IEEEauthorblockA{\textit{dept. name of organization (of Aff.)} \\
\textit{name of organization (of Aff.)}\\
City, Country \\
email address or ORCID}
\and
\IEEEauthorblockN{6\textsuperscript{th} Given Name Surname}
\IEEEauthorblockA{\textit{dept. name of organization (of Aff.)} \\
\textit{name of organization (of Aff.)}\\
City, Country \\
email address or ORCID}
}

\maketitle

\begin{abstract}
% Bối cảnh, động lực
% - Mạng camera biên đô thị và các nút IoT ngày càng phổ biến; yêu cầu phát hiện mục tiêu nhanh trong kịch bản khẩn cấp.
% - Quá trình truyền payload nhận diện chất lượng cao gặp giới hạn cửa sổ tiếp xúc ngắn giữa UAV và nút.
% 
% Bài toán và thách thức
% - Làm thế nào để nhanh chóng đưa dữ liệu hữu ích đến khu vực mục tiêu dưới ràng buộc PHY mức gói.
% - Thách thức: cửa sổ tiếp xúc ngắn, lỗi gói theo PER, 
% 
% Phương pháp đề xuất
% - Dữ liệu nhận diện được chia thành các fragment logic nhỏ hơn, mỗi fragment đều có giá trị nhận dạng nhất định.
% - UAV hoạt động như data-ferry, broadcast các fragment ra toàn khu vực
% - Nodes hợp tác thu thập và khôi phục dữ liệu 
% -> Phương pháp tổ chức mạng này đưa ra cơ chế phối hợp làm nhiệm vụ thay vì được giải quyết tập trung bởi UAV như những bài toán truyền thống.
% 
% Kết quả nổi bật
% - Phân tích mô phỏng cho thấy cơ chế phối hợp đóng vai trò quan trọng trong việc giảm thời gian hoàn thành nhiệm vụ, đặc biệt khi ràng buộc PHY nghiêm ngặt.
Urban edge-camera networks and IoT sensing infrastructures are increasingly deployed to support time-critical monitoring applications, where rapid target detection is essential in emergency scenarios. In such settings, Unmanned Aerial Vehicles (UAVs) can serve as mobile data ferries to disseminate high-value recognition payloads (e.g., images or feature representations used for target detection). However, transmitting large payloads is fundamentally constrained by short air-to-ground (A2G) contact windows and packet-level impairments, which limit reliable data delivery during brief flyovers.

This paper addresses the problem of minimizing mission-level time-to-detection under realistic PHY-layer constraints. The key challenge lies in delivering useful recognition information to ground nodes covering target areas despite limited contact duration and stochastic packet errors (PER), which make complete payload transfer impractical.

We propose a coordinated dissemination framework that jointly integrates UAV fragment broadcasting, coverage-aware trajectory planning, and cooperative edge fusion under packet-level feasibility constraints. Recognition payloads are partitioned into small fragments with approximately uniform information content, each contributing partial evidence toward detection. A UAV broadcasts these fragments along a coverage-aware trajectory, while ground nodes cooperatively exchange and fuse fragments to reconstruct missing information. Unlike approaches that rely primarily on UAV-only delivery, our design explicitly coordinates UAV dissemination and distributed edge cooperation under realistic packet-level constraints.

Simulation results show that the proposed coordination mechanism reduces time-to-detection, particularly under stringent PHY conditions with short contact windows and high packet error rates. The results highlight the importance of jointly optimizing dissemination, edge cooperation, and feasibility-aware communication in UAV-assisted urban sensing systems.

\end{abstract}

\begin{IEEEkeywords}
UAV-assisted dissemination, edge IoT, fragment-based recognition, cooperative detection
\end{IEEEkeywords}

\section{Introduction}

% Bối cảnh và bài toán
% - Đô thị thông minh ngày càng triển khai mạng camera IoT dày đặc phục vụ giám sát an ninh,
% trong đó việc phát hiện nhanh đối tượng khả nghi là yêu cầu cấp thiết. (hoặc có thể giới thiệu bài toán cứu hộ)
% - UAV nổi lên như một phương tiện cơ động có khả năng bao phủ khu vực rộng trong thời gian ngắn.
% - Tuy nhiên, payload nhận diện (ảnh/video chất lượng cao) có kích thước lớn,
% trong khi cửa sổ tiếp xúc giữa UAV và mỗi node chỉ kéo dài vài giây khi UAV bay qua.
% Nếu UAV dừng tại từng node để truyền đủ payload → mission time tăng mạnh, không đáp ứng yêu cầu realtime.
% Nếu UAV bay liên tục mà không dừng → mỗi node chỉ có thể nhận một phần payload trong cửa sổ ngắn đó.
% → Đây là trade-off cốt lõi mà bài toán cần giải quyết.

% Tại sao các hướng hiện tại chưa đủ tốt
% Các công trình UAV path planning (coverage, AoI, buffer-aware, service-time)
%   - tối ưu quỹ đạo tốt, nhưng thường mục tiêu ưu tiên rất rõ ràng (biết rõ mức độ ưu tiên của các node)
%   - thường giả định payload nguyên khối và mô hình kênh đơn giản (in-range / binary),
%   - được tính toán tập trung tại BS hoặc UAV, thiếu cơ chế phối hợp giữa các node để bù đắp cho ràng buộc PHY và tăng tốc độ phát hiện.
% Các công trình về hợp tác tại edge detection (threshold-triggered alert, peer collaboration, consensus)
%   - chưa xét đến overhead trao đổi hoặc giả định dữ liệu nhỏ 
% Các hướng đều bỏ qua ràng buộc PHY mức gói (A2G propagation, contact-window, SNR→BER→PER)
%   → dẫn đến đánh giá thuật toán không phản ánh đúng môi trường đô thị thực tế.

% Phương pháp đề xuất: 
% Giới thiệu cơ chế tích hợp các building block đã có dưới ràng buộc PHY thực tế để giải quyết bài toán nêu trên, bao gồm:
% - Tổ chức mạng mặt đất theo cell-layer 
% - Mô hình fragment-based recognition: chia payload nhận diện thành các fragment nhỏ hơn, mỗi fragment đều có giá trị nhận dạng nhất định (ví dụ: pixel-stride interleaving), giúp mỗi fragment đều có khả năng đóng góp vào confidence tích lũy.
% - UAV đóng vai trò "data ferry": bay theo quỹ đạo được tối ưu cho cell-layer, broadcast fragment liên tục khi di chuyển.

% Đóng góp chính
% Kết quả thử nghiệm cho thấy, cơ chế phối hợp này mang lại lợi ích đáng kể trong việc giảm thời gian hoàn thành nhiệm vụ. 
% Trong khi vẫn có thể giảm đáng kế quỹ đạo bay của UAV so với hướng không hợp tác
% đặc biệt khi ràng buộc PHY nghiêm ngặt (ví dụ: đường truyền kém, cửa sổ tiếp xúc ngắn). 
% Phân tích chi tiết về cách các tham số như kích thước fragment, ngưỡng hợp tác, và đặc tính kênh ảnh hưởng đến hiệu suất cũng được trình bày.

% P5 — Cấu trúc bài báo
% Section II: Related Work — 3 hướng và research gap.
% Section III: System Model — mạng, fragment, fusion, kênh PHY.
% Section IV: Problem Statement — bài toán tối ưu formal.
% Section V: Proposed Approach — GMC + cooperation protocol.
% Section VI: Performance Evaluation — setup, metrics, kết quả.
% Section VII: Discussion — giới hạn và hướng mở rộng.
% Section VIII: Conclusion.

Timely detection of critical events in urban environments increasingly relies on distributed sensing and edge-assisted inference. In such systems, spatially distributed ground nodes must rapidly acquire recognition payloads (e.g., images or feature representations) to identify targets of interest. This requirement is particularly prominent in applications such as security monitoring, search and rescue, and emergency response, where detection latency is critical.

A variety of approaches have been explored to address this challenge. Some works focus on performing inference directly at the edge to reduce communication latency. Others leverage distributed cooperation among nodes to improve detection accuracy. However, in complex urban environments, fixed communication infrastructures cannot always guarantee sufficient coverage or reliable connectivity, especially in emergency scenarios or rapid deployment settings.

In this context, Unmanned Aerial Vehicles (UAVs) have emerged as flexible mobile platforms to extend network coverage and act as data ferries to disseminate recognition payloads to ground nodes. While this approach helps overcome the limitations of static infrastructure, it also introduces a fundamental challenge: a mismatch between payload size and communication opportunity.

Air-to-ground (A2G) contacts between a UAV and ground nodes are typically brief—often lasting only a few seconds during flyovers—and are further affected by packet-level impairments such as path loss, shadowing, and fading. As a result, transmitting large payloads within a single contact window is often infeasible. Prolonging UAV dwell time increases mission latency, while partial transmissions leave individual nodes with insufficient evidence for reliable detection.

Our objective is to minimize mission-level time-to-detection under realistic A2G communication constraints.

These constraints suggest that attempting to deliver complete payloads to each node is neither efficient nor necessary. Instead, the key opportunity lies in exploiting partial information delivery and distributed cooperation.

Building on this insight, we develop a coordinated dissemination framework that integrates three complementary mechanisms under realistic PHY-layer constraints. First, recognition payloads are partitioned into multiple fragments with approximately uniform information content, allowing each fragment to contribute partial evidence toward detection. Second, the UAV continuously broadcasts fragments along a coverage-aware trajectory that balances spatial coverage and travel cost, avoiding the need for prolonged hovering. Third, ground nodes cooperatively exchange and fuse fragments, enabling them to reconstruct missing information through localized collaboration.

This coordinated design allows the system to combine UAV mobility with distributed edge cooperation, effectively compensating for short contact durations and stochastic packet losses. Nodes can progressively accumulate evidence and collaboratively improve detection confidence, enabling critical locations to reach decision thresholds more rapidly.

The core contribution of this work is a lightweight coordination layer that jointly schedules UAV fragment broadcasts, orchestrates cooperative exchange among ground nodes, and enforces packet-level feasibility based on contact duration and link conditions. By explicitly accounting for per-packet delivery constraints, the proposed approach avoids infeasible transmissions and limits unnecessary communication overhead.

Through extensive simulations, we show that the proposed framework significantly reduces mission-level time-to-detection compared to baseline strategies based on monolithic transmission or UAV-only delivery, particularly under stringent A2G conditions with short contact windows and high packet error rates.

Overall, this work advocates a shift from attempting complete payload delivery under stringent communication constraints to a design paradigm based on fine-grained dissemination and distributed cooperation, enabling timely and reliable detection in UAV-assisted urban sensing systems.

The remainder of the paper is organized as follows. Section~II reviews related work and positions our contribution. Section~III presents the system and channel models. Section~IV formulates the mission-time minimization problem. Section~V describes the proposed coordination protocol and trajectory design. Section~VI evaluates performance under realistic A2G conditions. Section~VII discusses limitations and extensions, and Section~VIII concludes.

\section{Related Work}
% Summarize related literature (skeleton only).
% Chia Related Work thành 3 hướng chính để bám sát bài toán hiện tại:
% - 1. UAV-assisted dissemination / data ferry / trajectory planning.
%   Điểm qua các nghiên cứu dùng UAV để mở rộng vùng phủ, thu thập hoặc phát dữ liệu trong mạng IoT/WSN,
%   đặc biệt là các bài tối ưu quỹ đạo theo độ phủ, độ trễ, service time, buffer-aware routing, hoặc time-balancing scheduling.
%   Ở đây chỉ nên nhắc AoI như một nhánh liên quan, không nên để AoI trở thành framing chính của bài.
%   Chốt ý: phần lớn các công trình này giả định payload nguyên khối hoặc mô hình liên lạc đơn giản,
%   chưa đi sâu vào nhiệm vụ phát hiện mục tiêu với dữ liệu lớn được phát theo từng fragment.
%
% - 2. Fragmented / progressive payload delivery và data fusion.
%   Điểm qua các nghiên cứu về chia nhỏ dữ liệu, progressive transmission, interleaving/subsampling,
%   rateless/fountain/network coding, hoặc partial reconstruction từ dữ liệu không đầy đủ.
%   Sau đó nối sang các cách hợp nhất thông tin từ fragment như union probability, Bayesian fusion, log-odds,
%   hoặc confidence accumulation để tạo quyết định từ bằng chứng từng phần.
%   Chốt ý: nhóm công trình này cho thấy dữ liệu phân mảnh vẫn có giá trị nhận dạng,
%   nhưng thường không xét UAV mobility, quỹ đạo phát tán, hay mục tiêu giảm time-to-first-detection.
%
% - 3. Distributed edge detection và cooperative inference.
%   Điểm qua các nghiên cứu về nhận dạng/phát hiện mục tiêu trên mạng cảm biến biên,
%   ngưỡng cảnh báo, multi-node cooperation, consensus, hoặc peer-assisted inference.
%   Nhấn mạnh rằng hợp tác giữa các node có thể cải thiện độ chính xác hoặc giảm false alarm,
%   nhưng đa số công trình giả định hạ tầng cố định và không gắn với cơ chế UAV broadcast fragment trong thời gian ngắn.
%
% Đoạn cuối của Related Work nên chốt thật rõ:
% - Các nghiên cứu về UAV path planning mạnh ở phía mobility, nhưng yếu ở mô hình payload phân mảnh và detection logic.
% - Các nghiên cứu về fragmented delivery / fusion mạnh ở phía dữ liệu và suy luận, nhưng thiếu ràng buộc UAV trajectory và A2G packet feasibility.
% - Các nghiên cứu về cooperative edge detection mạnh ở phía cộng tác giữa node, nhưng thường không tích hợp dissemination từ UAV và mô hình PHY mức gói.
% => Bài này nằm ở giao điểm của cả ba hướng: UAV phát fragment khi bay, node tích lũy confidence và hợp tác nội ô,
%    đồng thời tính đến ràng buộc PHY thực tế để tối ưu thời gian phát hiện đầu tiên.

% Prior studies on UAV-assisted IoT networking have explored coverage-aware trajectory design, service-time balancing, and latency-oriented routing. These methods provide strong mobility optimization tools, but most assume monolithic payload delivery or simplified communication success models, which limits their suitability for partial-information detection tasks.

% A second line of work investigates fragmented or progressive data delivery and evidence fusion. Such methods show that partial blocks can still contribute useful inference signal, and they often employ probabilistic aggregation rules to combine incomplete observations. However, these studies typically do not account for UAV mobility constraints, short contact windows, or mission-driven trajectory coupling.

% A third line of work addresses distributed edge detection and cooperative inference, including threshold-triggered alerts, peer collaboration, and consensus-style decision support. Although cooperation improves robustness and can reduce false alarms, most formulations assume static infrastructure and do not jointly model mobile aerial dissemination with packet-level air-to-ground feasibility.

% This paper lies at the intersection of these three directions. We jointly model UAV path planning, fragment-based confidence accumulation, and intra-cell cooperation under packet-level PHY constraints. The key focus is minimizing mission completion time defined by early confidence attainment at a designated suspicion point.

\section{System Model}
% Describe network/model assumptions and notation.
% 1. Thiết lập mạng: mạng được chia thành 2 phần: Ground node và UAV
% - Mô tả cách các edge node được triển khai dạng lưới đều trong khu vực đô thị 
% - các thông số liên quan đến vị trí và phạm vi hoạt động của UAV.
% 2. Ground node: 
% - Khu vực mạng được chia thành các cell, mỗi cell chứa một số nút biên. 
% - Mỗi nút có khả năng nhận diện đối tượng sử dụng fragments. 
% - Khu vực khả nghi được BS xác định trong đó có đối tượng bị truy nã (hoặc cần cứu hộ) nằm ở một node ngẫu nhiên (gọi là suspicion point).
% 3. Mô hình quỹ đạo UAV:
% - UAV có vùng phủ hữu hạn, bay liên tục và broadcast fragment trong khi di chuyển.
% - Quỹ đạo được tối ưu để cân bằng độ phủ vùng nghi vấn và chi phí di chuyển.
% 4. Mô hình dữ liệu: 
% - Lấy ví dụ về nhận diện đối tượng con người từ hình ảnh lớn có độ chính xác bao nhiêu
% - Kĩ thuật chia nhỏ dữ liệu: 
%     - pixel-stride interleaving để đảm bảo mỗi fragment có thông tin phân bố đều 
%     - giúp mỗi fragment đều có giá trị nhận dạng nhất định 
%     - nhưng lại có kích thước nhỏ hơn đáng kể và có khả năng fusion.
% 5. Mô hình hợp tác:
% - Điều kiện kích hoạt hợp tác: có 3 trường hợp:
%   - Khi node có độ tin cậy tích lũy vượt ngưỡng hợp tác, 
%   - Khi UAV vừa rời đi khỏi vùng phủ của node, hoặc
%   - Khi biết node khác có fragment mà mình thiếu (nhận được manifest).
% - Khi node trigger ngưỡng hợp tác, nó sẽ broadcast manifest fragment đang có (thông tin) cho các node trong cell để kích hoạt chia sẻ fragment.
% - Các node trong cell sẽ chia sẻ fragment với nhau theo yêu cầu
% 6. Mô hình kênh và khả năng nhận gói:
% - A2G/G2G profile với pathloss, shadowing, fast-fading.
% - Packet feasibility dựa trên contact-window và SNR $\rightarrow$ BER $\rightarrow$ PER.
% - Tốc độ truyền hiệu dụng là hàm theo khoảng cách, trạng thái kênh và chuyển động UAV theo thời gian.

We consider an urban IoT surveillance network composed of a base station (BS), one UAV, and ground-node set $\mathcal{N}$. Nodes are arranged on a grid and grouped into cells to enable local cooperation. The BS identifies a suspicious subset $\mathcal{P}\subseteq\mathcal{N}$ and designates one suspicion point $n^*\in\mathcal{P}$ as the mission-completion target.

\subsection{Network Setup and System Assumptions}
The monitoring field is partitioned into cells, each containing multiple edge nodes and an optional local coordinator for message scheduling. Each node can store fragment blocks, maintain a confidence score, and exchange blocks with peers in the same cell. The BS provides mission parameters (target set, thresholds, PHY settings, and UAV start state) before flight initialization.

\subsection{UAV Trajectory Planning Model}
The UAV flies at fixed altitude with finite broadcast radius $R_b$ and follows a waypoint sequence planned over suspicious nodes. Candidate waypoints include node positions and optional centroids. At planning step $t$, the next waypoint is selected by maximizing:
\begin{equation}
\text{score}(c)=\frac{\left|\mathrm{CS}(c)\setminus \mathrm{Covered}\right|}{\left(d(x_t,c)/v\right)^{\alpha}+\varepsilon},
\end{equation}
where $\mathrm{CS}(c)$ is the suspicious-node subset covered by candidate $c$, $d(x_t,c)$ is travel distance from the current UAV position, and $v$ is UAV speed. This favors candidates with high new coverage and low movement cost.

\subsection{Data and Fragment Dissemination Model}
The recognition payload is represented by $K$ logical fragments. The UAV broadcasts fragments continuously during flight rather than pausing for full-node transfer. Each node $n$ maintains received set $\mathcal{F}_n$ and updates confidence from observed fragments:
\begin{equation}
C_n = 1-\prod_{i\in\mathcal{F}_n}(1-p_i),
\end{equation}
where $p_i$ denotes the contribution of fragment $i$. When $C_n\ge\tau_{\text{coop}}$, node $n$ broadcasts a manifest and requests missing fragments from peers. Mission completion is triggered when the suspicion point reaches alert confidence, i.e., $C_{n^*}\ge\tau_{\text{alert}}$.

\subsection{Data Fusion Model}
To preserve broad spatial evidence in each block, fragments are modeled by pixel-stride interleaving over an image of size $W\times H$. Fragment $i$ contains pixels indexed by $j\equiv i\pmod{K}$, with size $s_i\approx WH/K$. This model approximates uniform information spread without requiring an end-to-end CV codec pipeline.

Each fragment contributes probability mass relative to its size:
\begin{equation}
p_i = 1 - (1 - C_{\text{master}})^{s_i/(WH)},
\end{equation}
where $C_{\text{master}}$ is confidence under complete payload availability. Hence, receiving all fragments recovers $C_n=C_{\text{master}}$.

Confidence fusion uses the union-probability rule:
\begin{equation}
C_n = 1 - \prod_{i \in \mathcal{F}_n}(1 - p_i).
\end{equation}
An equivalent log-odds representation under independence is:
\begin{equation}
C_n = \sigma\!\left(\sum_{i \in \mathcal{F}_n} \log \frac{p_i}{1-p_i}\right),
\end{equation}
where $\sigma(\cdot)$ is sigmoid. We use the union form in implementation for simplicity and numerical stability.

Node behavior is governed by two thresholds:
\begin{itemize}
  \item \textbf{Cooperation threshold} $\tau_{\text{coop}}$: trigger intra-cell manifest exchange and missing-fragment requests.
  \item \textbf{Alert threshold} $\tau_{\text{alert}}$: trigger alert; mission completes when this condition is met at $n^*$.
\end{itemize}
Setting $\tau_{\text{coop}}<\tau_{\text{alert}}$ enables two-stage amplification, reducing direct UAV delivery burden.

\subsection{Channel and Transmission-Rate Model}
The UAV--node link follows a geometry-aware air-to-ground model with path loss, shadowing, and fast fading. Received power is:
\begin{equation}
P_{rx}(t)=P_{tx}-PL(d_t)-X_{\sigma}(t)-X_f(t),
\end{equation}
where $d_t$ varies with UAV motion. SNR is mapped to BER and PER:
\begin{equation}
\mathrm{PER}=1-(1-\mathrm{BER})^{8L}.
\end{equation}
Packet acceptance additionally requires contact-window validity, i.e., receiver sensitivity must hold across packet airtime. Therefore, effective rate is time-varying with geometry and channel state.


\section{Problem Statement}
% Formal problem definition and objectives.
% - Mục tiêu: tối đa hóa thời gian phát hiện đầu tiên của đối tượng bị truy nã tại các nút biên trong khu vực khả nghi.
% - Ràng buộc: UAV có vùng phủ sóng giới hạn, dữ liệu nhận diện được chia thành các fragment
% - Bài toán: Lập kế hoạch quỹ đạo cho UAV để phát dữ liệu nhận và trao đổi dữ liệu giữa các nút biên nhằm đạt được độ tin cậy tích lũy vượt ngưỡng trong thời gian ngắn nhất có thể.
% - Điều kiện kết thúc: Khi node suspicious point đạt ngưỡng cảnh báo, nhiệm vụ được coi là hoàn thành.

Given suspicious set $\mathcal{P}$ and designated suspicion point $n^*\in\mathcal{P}$, the UAV broadcasts fragment set $\mathcal{F}=\{1,\ldots,K\}$ while moving. The objective is to minimize the time until $n^*$ reaches alert confidence. Decision variables are: (i) UAV trajectory over waypoints and (ii) fragment transmission schedule over time.

Let $\Pi=(w_1,w_2,\ldots,w_M)$ denote the UAV trajectory, where $w_m$ is the $m$-th waypoint, and let $\mathcal{S}=\{(f,\tau_f)\}$ denote the schedule transmitting fragment $f$ at time $\tau_f$. Let $C_n(t)$ be node confidence at time $t$, and $\tau_{\text{alert}}$ be alert threshold. Mission detection time is:
\begin{equation}
T_{\text{detect}}(\Pi,\mathcal{S}) = \inf\{t\ge 0\;|\; C_{n^*}(t)\ge \tau_{\text{alert}}\}.
\end{equation}

The optimization objective is to minimize expected detection time:
\begin{equation}
\min_{\Pi,\mathcal{S}}\;\mathbb{E}[T_{\text{detect}}(\Pi,\mathcal{S})].
\end{equation}

The problem is subject to the following operational constraints:
\begin{equation}
d(w_m,w_{m+1}) \le v_{\max}\,\Delta t_m,\;\forall m,
\end{equation}
\begin{equation}
\sum_{m=1}^{M-1}\frac{d(w_m,w_{m+1})}{v}+\sum_f t_f^{\text{tx}} \le T_{\max},
\end{equation}
\begin{equation}
\mathbf{1}_{\{n\leftarrow f\}}(t)=1 \Rightarrow P_{rx}^{(n)}(t')\ge P_{\text{sens}},\;\forall t'\in[t,t+t_f^{\text{tx}}],
\end{equation}
\begin{equation}
\Pr[\text{packet error}|n,f,t]=\mathrm{PER}(\mathrm{SNR}_{n,f}(t),L_f),
\end{equation}
\begin{equation}
\sum_{f\in\mathcal{F}_n(t)} b_f \le B_n^{\max},\;\forall n\in\mathcal{P},
\end{equation}
where $v_{\max}$ is maximum flight speed, $T_{\max}$ is mission deadline, $P_{\text{sens}}$ is receiver sensitivity, $L_f$ and $b_f$ are packet and fragment sizes, and $B_n^{\max}$ is node buffer capacity.

Because the problem is combinatorial (discrete trajectory, coupled scheduling, and stochastic channel constraints), global real-time optimization is intractable. We therefore design a lightweight heuristic that jointly approximates mobility and dissemination decisions.

\section{Proposed Approach}
\subsection{Overview}
% High-level description of the proposed method.
% (1) Input/Output
%     - Input: topology, suspicious set, suspicion point n^*, payload size, PHY parameters.
%     - Output: UAV flight path, confidence levels, mission completion time T_detect.
% (2) Pipeline tổng thể (end-to-end)
%     - BS tạo fragment
%     - UAV lập quỹ đạo GMC và broadcast fragment khi bay.
%     - Node cập nhật confidence + kích hoạt cooperation nội cell.
%     - Mission complete khi thỏa mãn điều kiện phát hiện.
% (3) Tại sao chọn heuristic thay vì tối ưu toàn cục
%     - Bài toán tổ hợp + ràng buộc kênh ngẫu nhiên + cần chạy realtime.
%     - GMC là xấp xỉ tốt, dễ mở rộng và deterministic theo seed.
% (4) Các module đóng góp chính
%     - Coverage-cost waypoint selection (GMC) + cell layer
%     - Fragment dissemination + confidence accumulation
%     - Intra-cell cooperation + packet-level feasibility gating

The proposed framework consists of three coordinated modules executed in a closed loop. First, the BS builds candidate waypoints and computes a GMC path for the UAV. Second, the UAV performs continuous fragment broadcasting along the path while each node updates confidence from successfully received packets. Third, nodes crossing the cooperation threshold exchange manifests and missing fragments inside each cell, accelerating confidence growth at the suspicion point.

Inputs include topology, suspicious-node set, suspicion point, threshold parameters, and PHY settings. Outputs include trajectory $\Pi$, transmission schedule $\mathcal{S}$, and mission detection time $T_{\text{detect}}$. We favor a heuristic design because exact optimization is computationally expensive under real-time mission constraints and random packet outcomes.

\subsection{Algorithm}
% Pseudocode / algorithm sketch.
% pseudo-code 2 mức
% - Algorithm 1: End-to-End Workflow)
% - đóng góp mới.
%
% A. Algorithm 1 -- End-to-End Workflow (ngắn, 12--20 dòng)
% Input: topology, suspicious set \mathcal{P}, suspicion point n^*, payload blocks, PHY params
% Output: \Pi, \mathcal{S}, T_detect
% Steps:
%   1) Initial Setup: tổ chức mạng ground node thành cells
%   2) Build candidates + precompute coverage
%   3) Data preparation: chia payload thành K fragment, mapping p_i, và C_master
%   4) Plan UAV path: heuristic GMC + cell-aware
%   5) Fragment dissemination: UAV broadcast theo schedule,
%   6) Tích lũy Confidence: 
%     loop: nhận fragment -> cập nhật confidence + kiểm tra: 
%                            - ngưỡng hợp tác -> broadcast manifest -> nhận fragment từ peer 
%                            - ngưỡng cảnh báo -> mission complete 
%
% B. Algorithm 2 -- Initial Setup
% - Ground nodes tự tổ chức thành cell layer dựa trên vị trí GPS
%   cell layer giúp tổ chức lập lịch giao tiếp tránh collision 
% - Mỗi node tự discover neighbors để bầu chọn cell leader
% - Fitness score cho CL được tính toán với tiêu chí phụ thuộc vào bài toán cụ thể 
%   (trong bài toán này thì dựa vào vị trí gần tâm cell nhất)
% - CL chịu trách nhiệm điều phối các hoạt động giao tiếp trong cell, 
%   và xác định gateway cho inter-cell cooperation.
%
% C. Algorithm 3 -- Path Planning
% - candidate set: node positions \cup optional k-means centroids
% - Loop: với mỗi candidate tính:
%   - coverage gain: với mỗi candidate, tính số node trong \mathcal{P} mà nó có thể phủ được (dựa trên bán kính R_b)
%   - Cell-aware: nếu một candidate có thê phủ được > 30% node trong một cell, thì đánh dấu toàn cell đã được phủ 
%   - greedy score: chọn ứng viên với gain/(cost^alpha + eps) tối ưu
% - output: ordered waypoints \Pi
%
% D. Algorithm 4 -- Cooperative Exchange
% Intra-cell cooperation được trigger khi:
% - C_n >= \tau_coop
% - UAV vừa rời khỏi vùng phủ của node 
% - Node nhận được manifest từ peer cho thấy node khác có fragment mà mình thiếu
% Khi trigger, node sẽ broadcast manifest (danh sách fragment đang có). 
% Các node trong cell sẽ chia sẻ fragment chưa có với nhau
% cell gateway (CGW) chủ động theo dõi tình hình của cả 2 cell thông qua neighbor cell gateway (NGW) 
% Inter-cell cooperation: CGW dựa vào thông tin có được quyết định và thực hiện khi:
% - C_n < \tau_coop && không còn manifes nào sinh ra trong cell
% - NGW cho thấy cell khác có fragment mà mình thiếu && C_n < \tau_alert && timeout 

\begin{algorithm}[t]
\caption{End-to-End Coordinated Detection Workflow}
\begin{algorithmic}[1]
\Require Topology $\mathcal{G}$, suspicious set $\mathcal{P}$, suspicion point $n^*$, payload $\mathcal{D}$, PHY parameters
\Ensure UAV path $\Pi$, schedule $\mathcal{S}$, detection time $T_{\text{detect}}$

\State \textbf{Initial Setup:} Organize ground nodes into cells $\{\mathcal{C}_i\}$ and elect cell leaders
\State \textbf{Precompute Coverage:} Build candidate set and coverage map over $\mathcal{P}$
\State \textbf{Data Preparation:} Partition $\mathcal{D}$ into $K$ fragments $\{f_i\}$ with uniform information content

\State \textbf{Path Planning:} Compute UAV trajectory $\Pi$ using GMC-based heuristic
\State \textbf{Initialize:} Set confidence $C_n \gets 0$ for all nodes

\For{each UAV waypoint $w \in \Pi$}
    \State Broadcast fragments according to schedule $\mathcal{S}$
    \For{each node $n$ receiving fragment $f_i$}
        \State Update confidence $C_n \gets C_n + \Delta(f_i)$

        \If{$C_n \ge \tau_{\text{coop}}$}
            \State Broadcast manifest and request missing fragments from peers
            \State Receive fragments and update $C_n$
        \EndIf

        \If{$n = n^*$ and $C_n \ge \tau_{\text{alert}}$}
            \State \Return $T_{\text{detect}}$
        \EndIf
    \EndFor
\EndFor
\end{algorithmic}
\end{algorithm}

\begin{algorithm}[t]
\caption{Cell Formation and Leader Election}
\begin{algorithmic}[1]
\Require Node set $\mathcal{N}$ with positions
\Ensure Cell structure $\{\mathcal{C}_i\}$ with leaders

\State Partition nodes into cells based on geographic proximity
\For{each cell $\mathcal{C}_i$}
    \State Nodes discover neighbors within communication range
    \State Compute fitness score for leader selection (e.g., proximity to cell centroid)
    \State Elect node with highest score as cell leader
    \State Designate gateway nodes for inter-cell communication
\EndFor
\end{algorithmic}
\end{algorithm}

\begin{algorithm}[t]
\caption{Coverage-Aware UAV Path Planning (GMC)}
\begin{algorithmic}[1]
\Require Candidate set $\mathcal{C}$, suspicious set $\mathcal{P}$
\Ensure Ordered waypoint list $\Pi$

\State Initialize $\Pi \gets \emptyset$, $\mathcal{P}_{\text{covered}} \gets \emptyset$

\While{$\mathcal{P}_{\text{covered}} \neq \mathcal{P}$}
    \For{each candidate $c \in \mathcal{C}$}
        \State Compute coverage gain $G(c)$ over uncovered nodes
        \If{$c$ covers $> \theta$ fraction of a cell}
            \State Mark entire cell as covered
        \EndIf
        \State Compute cost $L(c)$ (travel distance)
        \State Score $S(c) \gets \frac{G(c)}{L(c)^\alpha + \epsilon}$
    \EndFor

    \State Select $c^* = \arg\max S(c)$
    \State Append $c^*$ to $\Pi$
    \State Update $\mathcal{P}_{\text{covered}}$
\EndWhile
\end{algorithmic}
\end{algorithm}

\begin{algorithm}[t]
\caption{Cooperative Exchange (Intra- and Inter-cell)}
\begin{algorithmic}[1]
\Require Node $n$, fragment set $\mathcal{F}_n$, thresholds $\tau_{\text{coop}},\tau_{\text{alert}}$, cell gateway CGW, neighbor gateway NGW
\Ensure Updated fragment set(s) and confidence

\State define Trigger := ( $C_n \ge \tau_{\text{coop}}$ )
\State \quad \quad \quad OR (UAV just left $\Rightarrow$ contact window ended)
\State \quad \quad \quad OR (received manifest from peer showing missing fragments)

\If{Trigger}
    \State Broadcast manifest (list of owned fragments)
    \State Receive requests from neighbors in same cell
    \State Exchange missing fragments with peers (apply stagger/jitter scheduling to avoid collisions)
    \State Update local fragment set $\mathcal{F}_n$ and confidence $C_n$
\EndIf

\State /* Cell gateway (CGW) logic for inter-cell cooperation */
\If{CGW observes (no new manifests in cell) AND some nodes with $C_n < \tau_{\text{coop}}$}
    \If{NGW indicates neighbor cell has fragments missing locally AND target nodes satisfy $C_n < \tau_{\text{alert}}$ AND timeout expired}
        \State CGW request missing fragments from NGW
        \State NGW coordinates transfer (gateway-to-gateway exchange) and forwards fragments to requesting nodes
        \State Update fragment sets and confidences in both cells
    \EndIf
\EndIf
\end{algorithmic}
\end{algorithm}

\section{Performance Evaluation}
\subsection{Simulation Setup}
% Describe simulation environment and parameters.
% 1) Topology/network:
%    - số node, grid/cell layout, node spacing, seed, suspicion point selection policy.
% 2) UAV and mission:
%    - altitude, speed, broadcast radius, mission time budget, start position.
% 3) Communication/PHY:
%    - A2G/G2G profile parameters, sensitivity, contact-window config, BER/PER settings.
% 4) Data/fragment model (ghi rõ là abstract model):
%    - số fragment K, block size, confidence parameters (C_master, p_i mapping), thresholds \tau_coop/\tau_alert.
% 5) Baselines and variants:
%    - ví dụ NN path, random waypoint, GMC w/o cooperation, GMC full model.
%
% - số seeds, CI 95%, runtime budget mỗi scenario.
% - điều kiện dừng mô phỏng (mission complete hoặc timeout T_max).

\subsection{Metrics}
To evaluate the effectiveness of the proposed fragment-based cooperative detection and GMC planning, we employ the following metrics:
\begin{itemize}
    \item \textbf{Mission Completion Time ($T_{\text{complete}}$)}: The primary metric, defined as the time interval from UAV flight initialization until the designated suspicion point reaches the alert threshold $\tau_{\text{alert}}$.
    \item \textbf{UAV Path Length}: The cumulative distance traveled by the UAV along its waypoint sequence until the mission is completed.
    \item \textbf{Cooperation Gain}: The ratio of fragments received through intra-cell cooperative exchange versus those received directly from UAV broadcasts. This measures the effectiveness of the cell layer in compensating for link failures.
    \item \textbf{Packet Delivery Ratio (PDR)}: The fraction of successfully received fragment packets over the total number of packets broadcast by the UAV, illustrating the impact of mobility-aware PHY constraints.
\end{itemize}

\subsection{Results}

% kết quả sẽ có thể được so sánh với các baseline như: 
\begin{table}[t]
\centering
\caption{Mean mission completion time (seconds) by strategy and network size}
\label{tab:mission-times}
\small
\begin{tabular}{lrrrr}
\hline
Strategy & 100 & 400 & 900 & 1225 \\
\hline
Proposed (GMC w/ coop)        & 31.889  & 70.873   & 76.391   & 164.272 \\
NN (GreedyNearestNeighbor)     & 352.502 & 1331.598 & 2807.035 & 3431.988 \\
GMC w/o coop                   & 37.581  & 55.556   & 78.836   & 136.734 \\
Perfect channel (no packet loss) & 30.246  & 56.718   & 62.592   & 110.990 \\
\hline
\end{tabular}
\end{table}

\begin{table}[t]
\centering
\caption{Mean accumulated confidence (average) — With vs Without cooperation}
\label{tab:confidence-accum}
\small
\begin{tabular}{llccc}
\hline
Network size & Algorithm & After reception & At UAV depart & At mission complete \\
\hline
100  & With coop    & 0.441 & 0.477 & 0.802 \\
    & Without coop & 0.574 & 0.656 & 0.805 \\
\hline
400  & With coop    & 0.355 & 0.381 & 0.806 \\
    & Without coop & 0.584 & 0.623 & 0.805 \\
\hline
900  & With coop    & 0.387 & 0.415 & 0.804 \\
    & Without coop & 0.497 & 0.531 & 0.806 \\
\hline
1225 & With coop    & 0.379 & 0.416 & 0.805 \\
    & Without coop & 0.508 & 0.543 & 0.805 \\
\hline
\end{tabular}
\end{table}
\begin{table}[t]
\centering
\caption{Mean UAV flight path length (meters) by strategy and network size}
\label{tab:uav-path-lengths}
\small
\begin{tabular}{lrrrr}
\hline
Strategy & 100 & 400 & 900 & 1225 \\
\hline
Proposed (GMC w/ coop)        & 251.026  & 564.635   & 1120.984  & 1509.501 \\
NN (GreedyNearestNeighbor)     & 1105.195 & 3113.766 & 6344.382 & 8450.142 \\
GMC w/o coop                   & 485.808  & 1463.446 & 3016.069 & 3960.526 \\
Perfect channel (no packet loss) & 251.026  & 564.635   & 1120.984  & 1509.501 \\
\hline
\end{tabular}
\end{table}

% - brute-force nearest neighbor (NN) path planning
% - GMC không có hợp tác 
% - perfect channel (không có packet loss)

% - Sắp xếp hình/bảng theo ba phần: (A) kết quả cấp nhiệm vụ tổng quát, (B) phân tích theo lớp giao thức, (C) phân tích nhạy cảm / ablation.

% - (A) Thời gian hoàn thành nhiệm vụ ($T_{\text{complete}}$):cho cả 4 phương pháp (NN, GMC w/o coop, GMC full, perfect channel)
%     * Hình: CDF trung bình $T_{\text{complete}}$ 
%       - trục x: network size
%       - trục y: CDF của $T_{\text{complete}}$ 
%       - Dữ liệu: mean/median $T_{\text{complete}}$


% - (A2) Độ dài đường đi và chi phí: cho 3 phương pháp có tối ưu quỹ đạo (NN, GMC w/o coop, GMC full):
%     * Hình: CDF trung bình độ dài đường đi của UAV theo từng phương pháp chú thích mean/median.
%     - Dữ liệu: mean/median đường đi của UAV 


% - (B) Lợi ích hợp tác & phân phối fragment: so sánh giữa GMC w/o coop và GMC full 
%     * Stacked cluster bar + Time-series: 
%     - trung bình tích lũy độ tin cậy và tỉ lệ fragment confidence nhận từ UAV/hợp tác 
%     - trục x: thời điểm (hợp tác trước khi UAV đến, khi UAV rời đi, sau khi UAV rời đi, sau khi UAV hoàn thành quỹ đạo, tại $T_{\text{complete}}$)
%     - Dữ liệu: mean/median độ tin cậy tích lũy và tỉ lệ fragment confidence nhận từ UAV/hợp tác
%     * heatmap fairness
%     * pie chart phân phối trung bình fragment nhận được qua UAV vs cooperation của suspicion point $n^*$.

% - (B2) Các chỉ số mức gói (kiểm chứng PHY): (optional)
%     * PDR theo khoảng cách/độ cao: đường cong hoặc heatmap thể hiện Packet Delivery Ratio theo hình học UAV--node.
%     * Phân phối PER: histogram của PER trên các gói trong nhiều lần chạy; ghi chú mean/median.

% - (C) Phân tích nhạy cảm và ablation: (optional)
%     * Sweep plots: đồ thị đường của $T_{\text{complete}}$ (và tỉ lệ timeout) theo các tham số: $K$, $\tau_{\text{coop}}$, $\tau_{\text{alert}}$, $R_b$, mật độ nút.
%     * Bảng ablation: so sánh phương pháp đầy đủ với (i) không hợp tác, (ii) bỏ kiểm tra PER, (iii) tập ứng viên ngẫu nhiên; báo thay đổi tương đối % của $T_{\text{complete}}$.
% - Thống kê báo cáo:
%     * Sử dụng >=20 seed ngẫu nhiên cho mỗi kịch bản; tính 95% CI (bootstrap) cho các chỉ số chính.
%     * Báo p-value với kiểm định phù hợp (ví dụ Wilcoxon signed-rank) và đánh dấu trên hình khi p<0.05.

\section{Discussion}
% Limitations and interpretation.

\section{Conclusion}
% Summary and future work.

\section*{Acknowledgment}
% Acknowledgments (if any).

\bibliographystyle{IEEEtran}
\bibliography{references}

\appendix
\section{Appendix: Additional Material}
% Optional appendices.

\end{document}


% 1. Nền tảng về Giao tiếp UAV và Tối ưu hóa Năng lượng (Nhóm của Y. Zeng, R. Zhang) Đây là những công trình được trích dẫn nhiều nhất khi thiết lập mô hình năng lượng bay và tiềm năng mạng của UAV:
% "Wireless communications with unmanned aerial vehicles: Opportunities and challenges" (Y. Zeng, R. Zhang, T. J. Lim): Bài báo mang tính định hướng, chỉ ra cơ hội và thách thức của mạng giao tiếp bằng UAV
% .
% "Energy minimization for wireless communication with rotary-wing UAV" (Y. Zeng, J. Xu, R. Zhang): Nghiên cứu nền tảng về việc mô hình hóa công suất tiêu thụ cơ học và tối ưu hóa năng lượng cho UAV cánh quạt quay
% .
% "Energy-Efficient UAV Communication With Trajectory Optimization" (Y. Zeng, R. Zhang): Nghiên cứu tiên phong trong việc thiết kế quỹ đạo để đạt hiệu quả năng lượng cao nhất
% .
% 2. Mô hình Kênh truyền và Tối ưu hóa Độ cao bay (Nhóm của A. Al-Hourani) Hầu hết các bài báo khi giải quyết bài toán quỹ đạo 3D và kênh truyền Không gian - Mặt đất (Air-to-Ground) đều dựa vào các mô hình của nhóm tác giả này:
% "Optimal LAP altitude for maximum coverage" (A. Al-Hourani, S. Kandeepan, S. Lardner): Cung cấp công thức tính toán độ cao tối ưu cho các trạm phát sóng trên không tầm thấp (LAP) để đạt vùng phủ sóng tối đa
% .
% "Modeling air-to-ground path loss for low altitude platforms in urban environments" (A. Al-Hourani et al.): Đề xuất mô hình suy hao đường truyền (xác suất LoS/NLoS) nổi tiếng áp dụng cho môi trường đô thị
% .
% 3. Tuổi thông tin (AoI) trong Thu thập Dữ liệu (M. A. Abd-Elmagid, J. Liu, R. D. Yates) Với các bài nghiên cứu tối ưu hóa độ trễ và sự tươi mới của dữ liệu, các trích dẫn sau được coi là mốc chuẩn (benchmark):
% "Age of information: An introduction and survey" (R. D. Yates et al.): Bài tổng quan kinh điển định nghĩa và phân tích toán học về khái niệm Tuổi thông tin (AoI)
% .
% "Deep reinforcement learning for minimizing age-of-information in UAV-assisted networks" (M. A. Abd-Elmagid et al.) và "Age-optimal trajectory planning for UAV-assisted data collection" (J. Liu et al.): Các nghiên cứu đặt nền móng cho việc tối ưu quỹ đạo UAV dựa trên metric AoI thông qua Học tăng cường sâu (DRL)
% .
% 4. Nền tảng về Học tăng cường sâu - DRL (V. Mnih, H. V. Hasselt) Do rất nhiều bài sử dụng DRL để giải bài toán quỹ đạo, các bài báo gốc về thuật toán AI thường xuyên xuất hiện trong danh mục tài liệu tham khảo:
% "Human-level control through deep reinforcement learning" (V. Mnih et al., Nature 2015): Công trình đột phá của Google DeepMind giới thiệu mạng Deep Q-Network (DQN)
% .
% "Deep Reinforcement Learning with Double Q-learning" (H. V. Hasselt et al.): Bản nâng cấp của thuật toán DQN nhằm khắc phục nhược điểm đánh giá quá cao (overestimation) giá trị Q, được dùng rộng rãi trong điều hướng UAV
% .
% 5. Thiết kế phối hợp Quỹ đạo và Giao tiếp (Q. Wu, S. Jeong)
% "Joint trajectory and communication design for UAV-enabled multiple access/wireless networks" (Q. Wu, Y. Zeng, R. Zhang): Giải quyết bài toán phối hợp chặt chẽ giữa quỹ đạo di chuyển và lịch trình giao tiếp không dây
% .
% "Mobile edge computing via a UAV-mounted cloudlet: Optimization of bit allocation and path planning" (S. Jeong et al.): Đặt nền móng cho kiến trúc tích hợp tính toán biên (MEC) và thu thập dữ liệu lên UAV
% .

% Links — grouped by keyword

% UAV IoT survey:

% UAV-assisted data collection for Internet of Things: A survey — Z. Wei et al., IEEE (2022): https://ieeexplore.ieee.org/abstract/document/9779853/
% UAV IoT framework views and challenges — T. Lagkas et al., Sensors (2018): https://www.mdpi.com/1424-8220/18/11/4015
% AI for UAV-assisted IoT applications — N. Cheng et al. (2023) (IEEE / arXiv): https://ieeexplore.ieee.org/abstract/document/10113154/ and https://arxiv.org/pdf/2311.05303
% Secure communication in IoT-based UAV networks — J. Sharma et al., Internet of Things (2023): https://www.sciencedirect.com/science/article/pii/S2542660523002068
% UAV trajectory optimization (wireless) surveys:

% UAV placement and trajectory design optimization: A survey — A. Mazaherifar et al., Wireless Personal Communications (2022): https://link.springer.com/article/10.1007/s11277-021-09451-7
% A survey on UAV placement and trajectory optimization (A2G channel perspective) — J. Won et al., ICT Express (2023): https://www.sciencedirect.com/science/article/pii/S2405959522000157
% Survey on UAV deployment and trajectory in wireless networks — S.I. Han, Information (2022): https://www.mdpi.com/2078-2489/13/8/389
% Energy-efficient UAV communication with trajectory optimization — Y. Zeng & R. Zhang, IEEE Trans. Wireless Comm. (2017): https://ieeexplore.ieee.org/abstract/document/7888557

% UAV data-ferry / ferrying:

% A review of UAV ferry algorithms in delay tolerant network — B. Yang & X. Bai (2019): https://ieeexplore.ieee.org/abstract/document/9092564/
% Uav-assisted data collection in wireless sensor networks: A comprehensive survey — M.T. Nguyen et al., Electronics (2021): https://www.mdpi.com/2079-9292/10/21/2603
% End-to-end energy-efficiency and reliability of UAV-assisted wireless data ferrying — T. Shafique et al., IEEE (2019): https://ieeexplore.ieee.org/abstract/document/8930577/
% Opportunistic UAV utilization (systems overview) — D. Liu et al. (ResearchGate): https://www.researchgate.net/publication/342058665_Opportunistic-UAV-Utilization-in-Wireless-Networks_Motivations_Applications_and_Challenges
% Progressive transmission:

% Progressive image transmission: A review and comparison — K.H. Tzou, SPIE (1987): https://www.spiedigitallibrary.org/journals/optical-engineering/volume-26/issue-7/267581/Progressive-Image-Transmission-A-Review-And-Comparison-Of-Techniques/10.1117/12.7974121.short
% Survey of progressive image transmission methods — Y.K. Chee (1999): https://onlinelibrary.wiley.com/doi/abs/10.1002/(SICI)1098-1098(1999)10:1%3C3::AID-IMA2%3E3.0.CO;2-E
% The progressive smart grid system (IEEE Commun. Surveys & Tutorials, relevant progressive-transmission discussion): https://ieeexplore.ieee.org/abstract/document/5989902/

% Fragmented data / fragmentation surveys:

% A comprehensive taxonomy of fragmentation and allocation techniques — D. Nashat et al., ACM Computing Surveys (2018): https://dl.acm.org/doi/abs/10.1145/3150223
% Application of data fragmentation and replication methods in the cloud: a review — F. Castro-Medina et al., IEEE (2019): https://ieeexplore.ieee.org/abstract/document/8673249/
% (Related: fragmentation issues in IoT) Is Fragmentation a Threat to the Success of the Internet of Things? — M. Aly et al., IEEE Internet of Things (2018): https://ieeexplore.ieee.org/abstract/document/8424819/
% Fountain / rateless coding (wireless):

% Digital fountains: A survey and look forward — M. Mitzenmacher, IEEE (2004): https://ieeexplore.ieee.org/abstract/document/1405313/
% Fountain codes and their applications: Comparison and implementation for wireless applications — S. Chanayai & A. Apavatjrut, Wireless Personal Communications (2021): https://link.springer.com/article/10.1007/s11277-021-08749-w
% Primer and recent developments on fountain codes — J. Qureshi et al. (2013): https://www.ingentaconnect.com/content/ben/rptelec/2013/00000002/00000001/art00002
% Channel coding toward 6G: Technical overview and outlook (includes rateless/fountain discussion) — IEEE Open (2024): https://ieeexplore.ieee.org/abstract/document/10502324/

% Cooperative inference / edge surveys:

% DNN partitioning for cooperative inference in edge intelligence: Modeling, solutions, toolchains — ACM Computing Surveys (2026): https://dl.acm.org/doi/abs/10.1145/3786145
% Optimization methods, challenges, and opportunities for edge inference: A comprehensive survey — R. Zhang et al., Electronics (2025): https://www.mdpi.com/2079-9292/14/7/1345
% A survey on collaborative DNN inference for edge intelligence — W.Q. Ren et al., Machine Intelligence (2023): https://link.springer.com/article/10.1007/s11633-022-1391-7
% End-edge-cloud collaborative computing for deep learning: A comprehensive survey — Y. Wang et al., IEEE Commun. Surveys & Tutorials (2024): https://ieeexplore.ieee.org/abstract/document/10508191/
% CoEdge: Cooperative DNN inference with adaptive workload partitioning — IEEE/ACM Trans. (2020): https://ieeexplore.ieee.org/abstract/document/9296560
% Distributed detection / sensor-network surveys:

% Distributed sensor networks — a review of recent research — H. Qi et al., Journal of the Franklin Institute (2001) (classic): https://www.sciencedirect.com/science/article/pii/S0016003201000266
% Distributed detection with multiple sensors II: Advanced topics — R.S. Blum et al., Proc. IEEE (2002): https://ieeexplore.ieee.org/abstract/document/554209/
% Distributed event-triggered estimation over sensor networks: A survey — X. Ge et al., IEEE Transactions (2019): https://ieeexplore.ieee.org/abstract/document/8735960/
% A survey on distributed filtering and fault detection for sensor networks — H. Dong et al., Mathematical Problems in Eng. (2014): https://onlinelibrary.wiley.com/doi/abs/10.1155/2014/858624
% Application of compressive sensing techniques in distributed sensor networks (survey) — T. Wimalajeewa & P.K. Varshney, arXiv (2017): https://arxiv.org/abs/1709.10401

% Survey:
% "UAV-Assisted Data Collection for Internet of Things: A Survey" (Wei et al.)
% "A Comprehensive Survey on UAV-based Data Gathering Techniques in Wireless Sensor Networks" (Khan & Farooq)
% "A Review of UAV Ferry Algorithms in Delay Tolerant Network" (Yang & Bai)
% "AI for UAV-Assisted IoT Applications: A Comprehensive Review" (Cheng et al.)

% Tối ưu quỹ đạo theo độ phủ, độ trễ và service time:
% Trajectory optimization for the UAV assisted data collection in wireless sensor networks (Saxena et al.)
% EDGO: UAV-based effective data gathering scheme... (Pravija Raj et al.)
% Completion time minimization for multi-UAV-enabled data collection (Zhan & Zeng)
% End-to-End Energy-Efficiency and Reliability of UAV-Assisted Wireless Data Ferrying (Shafique et al.)

% Buffer-aware routing và Time-balancing scheduling:
% Research on a Storage Allocation Schema for DTN Ferry Node (Cong et al.)
% Energy minimization in Internet-of-Things system based on rotary-wing UAV (Zhan & Lai) 

% Tuổi của thông tin (AoI):
% Age-based path planning and data acquisition in UAV-assisted IoT networks (Jia et al.)
% UAV-enabled age-optimal data collection in wireless sensor networks (Tong et al.)




% - 2. Fragmented / progressive payload delivery và data fusion.
%   Điểm qua các nghiên cứu về chia nhỏ dữ liệu, progressive transmission, interleaving/subsampling,
%   rateless/fountain/network coding, hoặc partial reconstruction từ dữ liệu không đầy đủ.
%   Sau đó nối sang các cách hợp nhất thông tin từ fragment như union probability, Bayesian fusion, log-odds,
%   hoặc confidence accumulation để tạo quyết định từ bằng chứng từng phần.
%   Chốt ý: nhóm công trình này cho thấy dữ liệu phân mảnh vẫn có giá trị nhận dạng,
%   nhưng thường không xét UAV mobility, quỹ đạo phát tán, hay mục tiêu giảm time-to-first-detection.
%

% Hợp nhất thông tin từ các Fragment để tạo quyết định Để đưa ra quyết định từ các "bằng chứng từng phần" thay vì đợi toàn bộ dữ liệu, các phương pháp thống kê và tích lũy độ tự tin đã được đề xuất và khảo sát:
% Bayesian Fusion, Log-odds và Sequential Testing (Kiểm định tuần tự): Trong các mạng cảm biến phân tán, các nút biên hoặc trung tâm hợp nhất sẽ thu thập các quyết định/phân mảnh trung gian và sử dụng xác suất Bayes hoặc cập nhật Log-Likelihood Ratio (LLR) để ra quyết định dựa trên bằng chứng tích lũy theo thời gian
% "Distributed detection with multiple sensors II. Advanced topics"
% "Decentralized sequential detection with a fusion center performing the sequential test"
% "Bayesian compressive sensing via belief propagation"
% Confidence Accumulation (Tích lũy độ tự tin) và Early Exit: Trong suy luận phân tán của học sâu (Distributed Co-inference), mức độ tự tin từ các phân mảnh đặc trưng (feature maps) có thể được đo lường bằng hàm Entropy (mức độ không chắc chắn) hoặc hàm Softmax. Nếu độ tự tin vượt ngưỡng, hệ thống sẽ kết thúc sớm (early exit) quá trình suy luận. Các bài nghiên cứu nổi bật bao gồm:
% "BranchyNet: Fast inference via early exiting from deep neural networks"
% "SPINN: Synergistic progressive inference of neural networks over device and cloud"
% "DeeCap: Dynamic early exiting for efficient image captioning"
% "FrameExit: Conditional early exiting for efficient video recognition"


% - 3. Distributed edge detection và cooperative inference.
% . Hợp tác đa nút (Multi-node cooperation) và Đồng thuận phân tán (Consensus) trong Theo dõi/Phát hiện Mục tiêu Các hệ thống sử dụng bộ lọc đồng thuận (consensus filter) để các cảm biến phân tán trao đổi và hợp nhất các ước lượng trạng thái cục bộ, qua đó nâng cao khả năng quan sát toàn cục đối với mục tiêu di động:
% S. Zhu et al. đề xuất bộ lọc đồng thuận tối ưu phân tán cho theo dõi mục tiêu trong mạng cảm biến không đồng nhất (Distributed optimal consensus filter for target tracking in heterogeneous sensor networks)
% H. Su et al. nghiên cứu mạng cảm biến theo dõi đa mục tiêu bằng cơ chế lọc Kalman-consensus kích hoạt theo sự kiện (Event-triggered Kalman-consensus filter for two-target tracking sensor networks), giúp các cụm cảm biến vừa theo dõi chính xác vừa duy trì tần suất giao tiếp hợp lý
% Về phát hiện sự kiện, G. Wittenburg et al. nghiên cứu về phát hiện sự kiện hợp tác (Cooperative event detection) trong WSN
% , trong khi J. Meng et al. ứng dụng Compressive Sensing (cảm nhận nén) để phát hiện sự kiện thưa thớt từ các bằng chứng phân tán
% .
% 2. Ngưỡng cảnh báo (Alarm/Confidence Thresholds) và Giảm False Alarm Để giảm tỷ lệ cảnh báo sai (false alarm) hoặc đưa ra quyết định sớm khi có đủ bằng chứng, các hệ thống áp dụng các ngưỡng tin cậy tĩnh hoặc động:
% Cơ chế CFAR (Constant False Alarm Rate) phân tán: Trong nhận dạng tín hiệu radar/sonar, mạng đa cảm biến được dùng để thiết lập ngưỡng cảnh báo nhằm duy trì tỷ lệ false alarm không đổi trước nhiễu nền. Các công trình tiêu biểu bao gồm nghiên cứu của M. Barkat và P. K. Varshney về phát hiện tín hiệu CFAR phi tập trung
% ,
% , và M. K. Üner cùng P. K. Varshney về phát hiện CFAR phân tán trong các môi trường nhiễu đồng nhất và không đồng nhất
% .
% Early Exiting (Thoát sớm) trong Học sâu: Đối với suy luận trên thiết bị biên, các hệ thống đánh giá mức độ tin cậy của phân mảnh dữ liệu (dựa trên entropy hoặc Top-1 output). Nếu độ tự tin vượt ngưỡng, thiết bị sẽ ra quyết định ngay. Tiêu biểu là kiến trúc BranchyNet của Teerapittayanon et al.
%  hay hệ thống IoT học sâu theo yêu cầu EdgeKE của W. Fang et al. dùng trong môi trường dữ liệu công nghiệp lớn
% .
% 3. Peer-assisted Inference (Suy luận với thiết bị ngang hàng) Khái niệm mở rộng tài nguyên theo chiều ngang cho phép các thiết bị biên chia sẻ và thực thi song song các phân mảnh mô hình (model partitioning/chunks) hoặc dữ liệu để tăng tốc độ nhận dạng:
% Các hệ thống như DeepThings (Z. Zhao et al.)
% , MoDNN (J. Mao et al.)
% , và CoEdge (L. Zeng et al.)
%  cho phép phân chia linh hoạt khối lượng tính toán nội bộ. Các thiết bị lân cận (peer nodes) cùng chia sẻ tải tính toán của mạng CNN để phân tích các đặc trưng trích xuất từ dữ liệu cảm biến, cải thiện độ trễ tổng thể


% Đoạn cuối của Related Work nên chốt thật rõ:
% - Các nghiên cứu về UAV path planning mạnh ở phía mobility, nhưng yếu ở mô hình payload phân mảnh và detection logic.
% - Các nghiên cứu về fragmented delivery / fusion mạnh ở phía dữ liệu và suy luận, nhưng thiếu ràng buộc UAV trajectory và A2G packet feasibility.
% - Các nghiên cứu về cooperative edge detection mạnh ở phía cộng tác giữa node, nhưng thường không tích hợp dissemination từ UAV và mô hình PHY mức gói.
% => Bài này nằm ở giao điểm của cả ba hướng: UAV phát fragment khi bay, node tích lũy confidence và hợp tác nội ô,
%    đồng thời tính đến ràng buộc PHY thực tế để tối ưu thời gian phát hiện đầu tiên.
